\section{The Undertow Testbed}
\label{sec:system}

Answering C1--C4 requires an instrument, not an argument. Undertow is that
instrument: a testbed with two engines in which the same policy can be trained
cheaply and then measured against exact accounting on real on-chain events.
Figure~\ref{fig:system} sketches the data flow; this section describes each
stage and, just as importantly, the failure modes the testbed was built to
survive.

\begin{figure}[t]
\centering
\begin{tikzpicture}[
  font=\scriptsize,
  box/.style={draw, rounded corners=2pt, align=center, inner sep=3pt,
              text width=0.84\columnwidth},
  >={Latex[length=1.6mm]},
  node distance=3.2mm,
]
\node[box] (events) {on-chain events: swaps, mints, burns, collects};
\node[box, below=of events] (pipe) {data pipeline: event tape, reference prices, gas, regime labels, funding};
\node[box, below=of pipe] (view) {\texttt{MarketView}: the only door to history (look-ahead wall; train/eval split)};
\node[box, below=of view, fill=black!5] (engines) {exact Q128 backtester \quad\textbar\quad fast float64 Gymnasium simulator};
\node[box, below=of engines] (parity) {parity and look-ahead harness};
\node[box, below=of parity] (rl) {PPO policy over joint (timing, width, hedge ratio)};
\node[box, below=of rl] (eval) {evaluation: Sharpe, CVaR, cost decomposition, active/passive boundary};
\draw[->] (events) -- (pipe);
\draw[->] (pipe) -- (view);
\draw[->] (view) -- (engines);
\draw[->] (engines) -- (parity);
\draw[->] (parity) -- (rl);
\draw[->] (rl) -- (eval);
\end{tikzpicture}
\caption{Undertow: from on-chain events to a risk-sensitive policy. The two
engines share one \texttt{MarketView} and are reconciled by the parity harness
before any training result is reported.}
\label{fig:system}
\end{figure}

\subsection{Data: an on-chain event tape}

The data layer ingests the raw logs of two real USDC/WETH pools, the
$0.05\,\%$ primary study pool (\texttt{0x88e6a0\ldots f5640}, tick spacing
$10$) and the $0.30\,\%$ contrast pool (\texttt{0x8ad599\ldots e6d8}, spacing
$60$), and converts them into a single \emph{event tape} ordered by block and
containing swaps, mints, burns and collect events. Alongside the tape it
builds three auxiliary feeds: reference price closes at a resolution of one
minute, gas prices at each block, and regime labels produced by thresholds
committed in advance in the data module. A fourth feed, funding, is pulled from the Binance
perpetual market margined in USDT at its native $8\,\mathrm{h}$ cadence and
used by the hedge leg. The pinned window is \textbf{2022-01-01 to 2024-12-31
UTC}; the walk-forward split trains on 2022--2023 and evaluates on 2024, so the
test year is never seen during training or calibration. The primary pool's
training lane alone carries roughly $4.86\times 10^{6}$ events, which is why
the throughput work described below is a necessity rather than an optimization.

Two design rules protect the results. First, the tape is the ground truth:
prices, liquidity events and fees are the recorded on-chain values, not a
fitted process. Second, all access goes through \texttt{MarketView}, which
enforces both the block ordering and the train/eval boundary, so an
observation at time $t$ cannot read a bar from $t+1$ even by accident.

\subsection{The exact backtester}

The backtester replays a frozen policy against the tape using integer
arithmetic. Fee accrual is computed in Q128 fixed point from the pool's
recorded fee growth accumulators, exactly as the protocol does; liquidity
events update an external tick lattice seeded from a boundary snapshot, and
swap crossings are apportioned piecewise across the ticks they cross. The
resulting ledger satisfies an accounting identity by construction,
\begin{equation}
\texttt{net} = \texttt{fees} + \texttt{il} - \texttt{gas} - \texttt{slippage}
= \texttt{excess\_vs\_hodl},
\label{eq:ledger}
\end{equation}
with impermanent loss measured against the fixed initial deposit basket. The
policy protocol exposes no learning method, since a backtested policy is frozen
by construction, and every ledger records the config hash, seed and git commit
that produced it.

\subsection{The fast simulator}

Training on the exact engine is infeasible at millions of steps, so the
simulator mirrors the same economics in float64. Its agent is
\emph{marginal}: agent liquidity never moves the price path, and the fee share
is $\ell_{\text{agent}} / (\ell_{\text{agent}} + \ell_{\text{pool}})$ evaluated
against the tape's recorded active liquidity. The tick lattice supports both
minted agent positions and the external pool lattice, and the bookkeeping of
snapshots and the tape path uses exact Python integers so that a large gross
liquidity residue cannot silently corrupt a small delta, while a
burn larger than the position still raises. Two price modes share one
interface: \emph{replay}, which walks the reference closes at a resolution of
one minute, and a jump diffusion whose latent regimes follow a Markov chain
with four states anchored to the regime labels, with drift and volatility for
each state, Student-$t$ jumps, and transitions estimated from the training
split only. A throughput pass replaced
scans at each step with indexes computed as of that step, cached the lattice
keys and window views, and moved the vectorised environments into separate
processes; on the pinned hardware the PPO throughput of the full pipeline rose
from roughly $150$--$175$ to about $450$ steps per second. Actions execute at
the decision price observed \emph{before} the interval advances, matching the
backtester and the latency a live LP would face.

\subsection{Parity and look-ahead}

Two numeric worlds are only useful if their disagreement is bounded. The parity
harness feeds the simulator and the backtester the same price path, positions
and gas series and compares fees, impermanent loss and PnL under a stated
relative tolerance ($10^{-3}$ for float64 against exact integers). After the
fee engine was made exact, the remaining drift on the replay of real data is
small: mean fee differences within each seed are on the order of $10^{-7}$ of
notional, and impermanent loss matches exactly, comfortably inside tolerance.
The harness originally exercised only the HOLD policy, which never rebalances;
it now also runs a policy that rebalances, which is precisely the path the
C1--C3 fees depend on. Alongside parity, a suite of look-ahead probes attacks
the composed system: each component is checked to raise rather than read the
future, and the probes include deliberately adversarial windows at the
train/eval seam.

\subsection{Learning and the failure modes it survived}

The learning setup is deliberately modest: PPO with clipping on a $2\times256$
MLP, discount $0.99$, generalized advantage estimation (GAE) with
$\lambda = 0.95$, clip $0.15$, five fixed seeds, and one million timesteps per
seed. That simplicity is a feature: if a plain PPO cannot learn the task, the
task or the interface is broken, and it was. The first two runs collapsed
instead. Each seed locked onto one deterministic action, some always holding
and some always repeating a single rebalance, and the measured advantage did
not separate the actions, with a correlation across actions near zero.
Diagnosing that collapse produced three interface fixes that are now part of
the testbed and are ablated for C2. The observation gained the three band
geometry features, and the features with heavy tails gained signed log
compression. The reward term for impermanent loss was moved from a rolling
entry basket to the fixed initial HODL basket so that the reward at each step
equals the change in the reported excess metric; the old parity harness could not
catch that divergence because it only ran HOLD. Episodes gained window
resampling: a reset without an episode seed now advances to a new window of
30 days instead of replaying the same lane for a million steps. Two cost defects were
fixed on the way: rebalance slippage is charged on the inventory imbalance
actually swapped rather than the whole position value, and the marginal fee
credit is capped at the augmented denominator share so that a tiny agent
cannot claim an outsized fee fraction. Each fix ships with a unit test and a
decision record; the C2 table in Section~\ref{sec:results} lists symptom,
diagnosis, fix and evidence.

\subsection{Evaluation and reproducibility}

The evaluation runner freezes a checkpoint and scores it on windows held out
from training, emitting each episode's equity curve, cost ledgers and regime
slices, then aggregating annualized Sharpe, raw PnL and excess-versus-HODL with
confidence intervals across seeds and regimes. \pending{CVaR at $5\,\%$ exists
as a metric in the metrics layer; wiring it into the aggregated evaluation
report is the S21 deliverable, after which it joins the primary statistics.}
The RQ1 ablation table, the RQ2 regime matrix and the report on the gap
between simulator and backtester are generated from the same runner, and a
dashboard that is read only replays policies against price and ranks runs by
Sharpe out of sample. Every artefact carries the config hash and commit, so a
reported number can be traced to the exact run that produced it.
